\documentclass{article}
\usepackage{amsmath, amssymb, amsthm}
\usepackage{hyperref}
\usepackage{geometry}
\geometry{margin=1in}

\title{Erd\H{o}s Problem \#609}
\date{Last updated: 2025-08-31}
\author{Source: erdosproblems.com}

\begin{document}
\maketitle

\noindent\textbf{Status:} OPEN \quad \textbf{No prize}

\noindent\textbf{Tags:} graph theory, ramsey theory

\medskip
\noindent\textbf{Problem Statement:}

Let $f(n)$ be the minimal $m$ such that if the edges of $K_{2^n+1}$ are coloured with $n$ colours then there must be a monochromatic odd cycle of length at most $m$. Estimate $f(n)$.

\bigskip
\noindent\textbf{Remarks:}

A problem of Erd\H{o}s and Graham. The edges of $K_{2^n}$ can be $n$-coloured to avoid odd cycles of any length. It can be shown that $C_5$ and $C_7$ can be avoided for large $n$.

Chung \cite{Ch97} asked whether $f(n)\to \infty$ as $n\to \infty$. Day and Johnson \cite{DaJo17} proved this is true, and that\[f(n)\geq 2^{c\sqrt{\log n}}\]for some constant $c>0$. The trivial upper bound is $2^n$.

Gir\~{a}o and Hunter \cite{GiHu24} have proved that\[f(n) \ll \frac{2^n}{n^{1-o(1)}}.\]Janzer and Yip \cite{JaYi25} have improved this to\[f(n) \ll n^{3/2}2^{n/2}.\]See also the entry in the graphs problem collection.

\bigskip
\noindent\textbf{References:}

[Ch97] Chung, F. R. K., Open problems of {P}aul Erd\H{o}s in graph theory. J. Graph Theory (1997), 3--36.
 
 [DaJo17] Day, A. Nicholas and Johnson, J. Robert, Multicolour {R}amsey numbers of odd cycles. J. Combin. Theory Ser. B (2017), 56--63.
 
 [GiHu24] A. Gir\~Ao and Z. Hunter, Monochromatic odd cycles in edge-coloured complete graphs. arXiv:2412.07708 (2024).
 
 [JaYi25] O. Janzer and F. Yip, Short monochromatic odd cycles. arXiv:2506.14910 (2025).

\bigskip
\noindent\small{Source: \url{https://www.erdosproblems.com/609}}
\end{document}
